\documentclass{umk-matematika}

\begin{document}

\maketitlepage
    {Практическая работа №12}
    {Перпендикулярность прямой и плоскости}
    {}

\section{Фундамент (1--4 балла)}

\textbf{Задача 1 (1 балл).} Сформулируйте определение перпендикулярности прямой и плоскости.

\textbf{Задача 2 (2 балла).} Сформулируйте признак перпендикулярности прямой и плоскости.

\textbf{Задача 3 (3 балла).} Что такое перпендикуляр, опущенный из точки на плоскость?

\textbf{Задача 4 (4 балла).} Что такое наклонная к плоскости? Что такое проекция наклонной?

\section{Стены (5--7 баллов)}

\textbf{Задача 5 (5 баллов).} Из точки A к плоскости α проведены перпендикуляр AB = 6 см и наклонная AC = 10 см. Найдите длину проекции наклонной BC.

\textbf{Задача 6 (6 баллов).} Прямая a перпендикулярна плоскости α. Что можно сказать о взаимном расположении прямой a и любой прямой, лежащей в плоскости α?

\textbf{Задача 7 (7 баллов).} В кубе $ABCDA_1B_1C_1D_1$ докажите, что ребро AA₁ перпендикулярно плоскости ABCD.

\section{Кровля (8--10 баллов)}

\textbf{Задача 8 (8 баллов).} Из точки M к плоскости проведены две наклонные MA = 13 см и MB = 15 см. Их проекции относятся как 5:9. Найдите расстояние от точки M до плоскости.

\textbf{Задача 9 (9 баллов).} Через вершину A прямоугольника ABCD проведён перпендикуляр AK к плоскости прямоугольника. Найдите расстояние от точки K до вершины C, если AB = 6 см, AD = 8 см, AK = 12 см.

\textbf{Задача 10 (10 баллов).} Строительный отвес длиной 2 м отклонили от вертикали так, что его верхний конец сместился на 30 см от стены. На каком расстоянии от стены находится нижний конец отвеса?

\newpage
\section*{Ответы}

\begin{enumerate}
  \item определение перпендикулярности.
  \item признак перпендикулярности.
  \item кратчайшее расстояние от точки до плоскости.
  \item наклонная длиннее перпендикуляра; проекция лежит в плоскости.
  \item $BC = 8$ см.
  \item прямая a перпендикулярна любой прямой в плоскости α.
  \item AA₁ ⟂ (ABCD).
  \item $MO = 12$ см.
  \item $KC = 2\sqrt{61}$ см.
  \item 30 см.
\end{enumerate}

\end{document}